% ============================================================
%  朱赫 — 个人简历 (中文版)
%  Style: clean academic CV (matching English cv.tex layout)
% ============================================================
\documentclass[10pt, a4paper]{article}

% ── Encoding & CJK Support ──────────────────────────────────
\usepackage[UTF8, heading=false]{ctex}
\usepackage[T1]{fontenc}
\usepackage{palatino}
\usepackage{mathpazo}
\usepackage{microtype}

% ── Page Layout ──────────────────────────────────────────────
% Uniform tight margins so page 2 can fit all publications + honors
\usepackage[top=0.7cm, bottom=0.7cm, left=1.5cm, right=1.5cm]{geometry}
\setlength{\parindent}{0pt}
\setlength{\parskip}{0pt}

% ── Packages ─────────────────────────────────────────────────
\usepackage{xcolor}
\usepackage{hyperref}
\usepackage{enumitem}
\usepackage{titlesec}
\usepackage{tabularx}
\usepackage{array}
\usepackage{fancyhdr}
\usepackage{tikz}
\usepackage{etoolbox}
\usepackage{graphicx}

% ── Colors ───────────────────────────────────────────────────
\definecolor{headcol}{RGB}{30, 30, 30}
\definecolor{rulecol}{RGB}{180, 50, 30}
\definecolor{secline}{RGB}{200, 200, 200}
\definecolor{linkcol}{RGB}{40, 80, 160}
\definecolor{dategray}{RGB}{100, 100, 100}
\definecolor{yearcolor}{RGB}{80, 80, 80}

% Tag colors
\definecolor{tagalign}{RGB}{200, 60, 30}
\definecolor{tagdata}{RGB}{30, 110, 180}
\definecolor{tagurban}{RGB}{40, 140, 90}
\definecolor{tageff}{RGB}{130, 50, 160}
\definecolor{tagpost}{RGB}{180, 100, 20}
\definecolor{oralred}{RGB}{200, 40, 30}

\hypersetup{
  colorlinks = true,
  urlcolor   = linkcol,
  linkcolor  = linkcol,
  pdfauthor  = {He Zhu},
  pdftitle   = {朱赫 — 个人简历},
}

% ── Section Headers ──────────────────────────────────────────
\titleformat{\section}
  {\normalsize\bfseries\color{headcol}}
  {}{0em}{}
  [\vspace{2pt}{\color{secline}\rule{\linewidth}{0.4pt}}\vspace{-2pt}]
\titlespacing*{\section}{0pt}{6pt}{3pt}

% ── Badge Tags ───────────────────────────────────────────────
\newcommand{\badge}[2]{%
  \tikz[baseline=(x.base)]{
    \node[inner xsep=3.5pt, inner ysep=1.6pt, rounded corners=2.5pt,
          draw=#1, line width=0.45pt,
          text=#1, font=\fontsize{6}{7}\selectfont\sffamily]
      (x) {#2};
  }%
}

\newcommand{\talign}{\badge{tagalign}{alignment}}
\newcommand{\tdata} {\badge{tagdata}{data-centric}}
\newcommand{\turban}{\badge{tagurban}{urban AI}}
\newcommand{\teff}  {\badge{tageff}{efficiency}}
\newcommand{\tpost} {\badge{tagpost}{post-train}}

% ── Bullet Lists ─────────────────────────────────────────────
\setlist[itemize,1]{
  leftmargin = 1.3em,
  label      = {\small\textbullet},
  itemsep    = 1pt,
  topsep     = 2pt,
  parsep     = 0pt,
}

% ── Helpers ──────────────────────────────────────────────────
\newcommand{\unilogo}[1]{%
  \raisebox{-0.5pt}{\includegraphics[height=8pt]{assets/logos/cv/#1}}%
}
\newcommand{\me}[1]{\textbf{#1}}
\newcommand{\eqstar}{$^{*}$}
\newcommand{\corr}{$^{\dagger}$}

% Configurable publication vertical skip
\def\pubskip{3pt}

% Single publication block
\newcommand{\pub}[4]{%
  \vspace{\pubskip}%
  \noindent\textbf{#1}\ifstrempty{#2}{}{\enskip #2}%
  \begin{itemize}[topsep=0pt]
    \item #3
    \item \textit{#4}
  \end{itemize}%
}

% Co-authored publication (no author line)
\newcommand{\copub}[3]{%
  \vspace{\pubskip}%
  \noindent\textbf{#1}\ifstrempty{#2}{}{\enskip #2}%
  \begin{itemize}[topsep=0pt]
    \item \textit{#3}
  \end{itemize}%
}

% Education / Experience entry
\newcommand{\entry}[2]{%
  \noindent
  \begin{tabularx}{\linewidth}{@{}X r@{}}
    \textbf{#1} & {\small\color{dategray}#2}
  \end{tabularx}%
  \par\vspace{0pt}%
}

% ── Header / Footer ──────────────────────────────────────────
\pagestyle{fancy}
\fancyhf{}
\renewcommand{\headrulewidth}{0pt}
\fancyfoot[R]{\footnotesize\color{dategray}\thepage}

% ============================================================
\begin{document}
\thispagestyle{empty}

% ── NAME BLOCK ───────────────────────────────────────────────
\noindent
\begin{tabularx}{\linewidth}{@{}X r@{}}
  \raisebox{-4pt}{%
    {\fontsize{28}{32}\selectfont\bfseries 朱赫}%
    \quad{\large\color{dategray}2027应届}%
  }
  &
  \begin{tabular}[b]{@{}r@{}}
    \small\href{mailto:zhuhe@stu.pku.edu.cn}{zhuhe@stu.pku.edu.cn}\\[2pt]
    \small\href{https://zhuchichi56.github.io}{https://zhuchichi56.github.io}
  \end{tabular}
\end{tabularx}

\vspace{6pt}
{\color{rulecol}\rule{\linewidth}{1.2pt}}
\vspace{4pt}

% ── ABOUT ────────────────────────────────────────────────────
\section{个人简介}

24岁，\textbf{北京大学}硕士研究生（导师：张文佳教授、陈冠华教授），
本科毕业于\textbf{南方科技大学}计算机科学与工程系，
现于\textbf{微软亚洲研究院} GenAI 组（Dr.~Li Dong）担任研究实习生。
累计发表大模型算法相关论文第一作者/通讯作者 9 篇（含 2 篇 under review）。
曾获腾讯青云实习生计划（CSIG, 2025）、美团北斗实习生计划（M17, 2026）、
九坤基础语言模型实习生（2026）等 offer。

\vspace{3pt}
研究聚焦于\textbf{大模型后训练（LLM Post-Training）算法}，涵盖两条主线：
\textbf{(1) 高质量后训练数据的合成与筛选}——如何定义高质量的后训练数据，
并提出可规模化的自动合成与质量筛选方法
(\href{https://aclanthology.org/2025.findings-acl.906/}{FANNO},
 \href{https://arxiv.org/abs/2601.23006}{InstructDiff},
 \href{https://aclanthology.org/2025.findings-acl.911/}{Tag-Instruct},
 \href{https://openreview.net/forum?id=4fclVIrUg2}{AlignDiff})；
\textbf{(2) 高效训练目标设计}——探索如何在习得新能力的同时缓解灾难性遗忘
(\href{https://arxiv.org/abs/2509.23753}{ASFT},
 \href{https://arxiv.org/abs/2510.10974}{CFT})。
此外，将上述后训练研究范式延伸至\textbf{垂直领域大模型}建设，
在城市智能场景中完成了从数据构建、模型训练到系统落地的全链路实践
(\href{https://arxiv.org/abs/2402.19273}{PlanGPT},
 \href{https://arxiv.org/abs/2505.14481}{PlanGPT-VL},
 \href{https://app.urbanclaw.net}{UrbanClaw})。

% ── EDUCATION ────────────────────────────────────────────────
\section{教育经历}

\entry{\unilogo{pku.png}\enskip 北京大学}{北京}
\textit{理学硕士，智慧城市与数据科学} \hfill \textit{2024 -- 至今}
\begin{itemize}
  \item \textbf{导师}：张文佳教授、陈冠华教授
  \item \textbf{研究方向}：Data-centric LLM post-training, on-policy distillation, alignment
\end{itemize}

\vspace{4pt}
\entry{\unilogo{sustech.png}\enskip 南方科技大学}{深圳}
\textit{工学学士，计算机科学与工程} \hfill \textit{2020 -- 2024}
\begin{itemize}
  \item \textbf{导师}：范子沛教授、宋轩教授
  \item GPA 90.2/100（前10\%）；优秀毕业生、优秀毕业设计（前5\%，2024）
  \item \textbf{访问}：\unilogo{utokyo.png}\,\href{https://www.csis.u-tokyo.ac.jp/}{东京大学 CSIS}（2022, 2023）；\unilogo{nus.png}\,\href{https://www.comp.nus.edu.sg/}{新加坡国立大学 CS}（2022）
\end{itemize}

% ── EXPERIENCE ───────────────────────────────────────────────
\section{实习与工作经历}

\entry{\unilogo{microsoft.png}\enskip 微软亚洲研究院 $\cdot$ GenAI 组}{北京}
\textit{研究实习生}，导师：Dr.~Li Dong \hfill \textit{2025.11 -- 至今}\\[1pt]
On-policy distillation \& agent harness

\vspace{3pt}
\entry{\unilogo{pjlab.png}\enskip 上海人工智能实验室 $\cdot$ OpenData Lab}{上海}
\textit{研究实习生}，foundation language models \hfill \textit{2025.05 -- 2025.10}

\vspace{3pt}
\entry{\unilogo{sensetime.png}\enskip 商汤科技 $\cdot$ Foundation Language Model Center}{深圳}
\textit{研究实习生}，foundation models \hfill \textit{2024.06 -- 2024.09}

\vspace{3pt}
\entry{\unilogo{locationmind.png}\enskip LocationMind}{东京，日本}
\textit{项目助理}，autonomous driving algorithms \hfill \textit{2023.06 -- 2023.12}

% ── PROJECTS ─────────────────────────────────────────────────
\section{代表项目}

\noindent\textbf{PlanGPT 系列} \hfill
{\small\href{https://plangpt.github.io}{plangpt.github.io}}\\[2pt]
首个系统性研究 LLM 赋能城市规划的工作。包含
\textbf{PlanGPT}（面向规划方案生成的文本大模型）、
\textbf{PlanGPT-VL}（面向规划图件理解的 vision-language model）、
\textbf{PlanGPT-R1}（reasoning-enhanced model）及
\textbf{UP-Bench / PlanBench-V}（benchmarks）。

\vspace{3pt}
\noindent\textbf{UrbanClaw} \hfill
{\small\href{https://app.urbanclaw.net}{app.urbanclaw.net}}\\[2pt]
AI 驱动的城市规划 agent，具备 multi-agent collaboration、Docker sandbox、
MCP client layer 及 self-evolving skill system。

% ============================================================
% PAGE 2 — switch to tighter spacing
% ============================================================
\newpage
\def\pubskip{1pt}
\titlespacing*{\section}{0pt}{6pt}{2pt}

% ── PUBLICATIONS ─────────────────────────────────────────────
\section{学术论文}

\noindent\textit{第一作者 / 通讯作者}

\pub{PlanGPT: Enhancing Urban Planning with a Tailored Language Model}{\turban}
    {\me{He Zhu}, Guanhua Chen, Wenjia Zhang\corr}
    {ACL Industry 2025 $\cdot$ {\bfseries\color{oralred}Oral} \quad[\href{https://arxiv.org/abs/2402.19273}{paper}]}

\pub{Anchored Supervised Fine-Tuning}{\talign}
    {\me{He Zhu}\eqstar, Junyou Su\eqstar, Peng Lai\eqstar,
     Wenjia Zhang, Linyi Yang, Guanhua Chen\corr}
    {ICLR 2026 \quad[\href{https://arxiv.org/abs/2509.23753}{paper}]}

\pub{InstructDiff: Domain-Adaptive Data Selection via Differential Entropy}{\tdata}
    {Junyou Su\eqstar, \me{He Zhu}\eqstar\corr, Guanhua Chen\corr}
    {ACL 2026 (Main) \quad[\href{https://arxiv.org/abs/2601.23006}{paper}]}

\pub{FANNO: Augmenting High-Quality Instruction Data with Open-Sourced LLMs Only}{\tdata}
    {\me{He Zhu}, Yifan Ding, Yicheng Tao, \ldots, Guanhua Chen\corr}
    {ACL 2025 (Findings) \quad[\href{https://aclanthology.org/2025.findings-acl.906/}{paper}]}

\pub{Tag-Instruct: Controlled Instruction Complexity Enhancement}{\tdata}
    {\me{He Zhu}, Zhiwen Ruan, Junyou Su, \ldots, Guanhua Chen\corr}
    {ACL 2025 (Findings) \quad[\href{https://aclanthology.org/2025.findings-acl.911/}{paper}]}

\pub{PlanGPT-VL: Vision-Language Model for Urban Planning Maps}{\turban}
    {\me{He Zhu}\eqstar, Junyou Su\eqstar, Minxin Chen\eqstar,
     \ldots, Wenjia Zhang\corr}
    {EMNLP Industry 2025 \quad[\href{https://arxiv.org/abs/2505.14481}{paper}]}

\pub{Personalized Individual Trajectory Prediction via Meta-Learning}{\turban}
    {\me{He Zhu}, Liyu Zhang, Zipei Fan}
    {SIGSPATIAL 2022 $\cdot$ {\bfseries\color{oralred}Oral} \quad[\href{https://dl.acm.org/doi/abs/10.1145/3557915.3565536}{paper}]}

\pub{AlignDiff: Exploiting Model-Intrinsic Information for Better Data Selection}{\tdata}
    {Peng Lai\eqstar, \me{He Zhu}\eqstar, Zhiwen Ruan, \ldots, Yang Liu, Guanhua Chen\corr}
    {Under Review, 2026 \quad[\href{https://openreview.net/forum?id=4fclVIrUg2}{paper}]}

\pub{Towards Fair and Comprehensive Evaluation of Routers
     in Collaborative LLM Systems}{\teff}
    {Wanxing Wu\eqstar, \me{He Zhu}\eqstar, Yixia Li\eqstar,
     Yun Chen, Guanhua Chen\corr}
    {Under Review, 2026 \quad[\href{https://arxiv.org/abs/2602.11877}{paper}]}

\vspace{2pt}
\noindent\textit{合作论文（精选）}

\copub{Enhancing LLM Reasoning via Selective Critical Token Fine-Tuning}{\talign}
    {Under Review, 2026 \quad[\href{https://arxiv.org/abs/2510.10974}{paper}]}

\copub{Dripper: Token-Efficient Main HTML Extraction with a Lightweight LM}{\talign}
    {Under Review, 2026 \quad[\href{https://arxiv.org/abs/2511.23119}{paper}]}

\copub{Topic Over Source: The Key to Effective Data Mixing for LLM Pre-training}{\tdata}
    {Under Review, 2026 \quad[\href{https://arxiv.org/abs/2502.16802}{paper}]}

\copub{LayAlign: Enhancing Multilingual Reasoning via Layer-wise Adaptive Fusion}{\talign}
    {NAACL 2025 (Findings)}

\copub{ToolExpNet: Optimizing Multi-Tool Selection in LLMs with Experience Networks}{\teff}
    {ACL 2025 (Findings)}

\copub{HHGNN: Heterogeneous Hypergraph Neural Network for Trajectory Prediction}{\turban}
    {ICRA 2024}

% ── HONORS & SERVICE ─────────────────────────────────────────
\section{荣誉与学术服务}

\begin{itemize}[itemsep=0pt, topsep=1pt]
  \item 北京大学优秀团支书（\textbf{前5\%}，2025）；北京大学学生代表（2025）
  \item 南方科技大学优秀毕业生、优秀毕业设计（\textbf{前5\%}，2024）；年度优秀学生（2021, 2022, 2023）
  \item \textbf{审稿人}：ACL Rolling Review (ARR, 2024--2026)
\end{itemize}

\end{document}
